% Sections/05-KhaiThacYeuCau.tex
\section{Khai thác yêu cầu}

Chương này xác định tác nhân, yêu cầu chức năng, phi chức năng và dữ liệu của hai hệ thống Wispace Speaking và AIHub. Wispace phục vụ luyện tập, thi thử và đánh giá kỹ năng Nói IELTS; AIHub cung cấp API mô hình cho đối tác B2B và công cụ quản lý chất lượng mô hình nội bộ. Các yêu cầu làm cơ sở cho phân tích nghiệp vụ ở Chương 6 và thiết kế ở Chương 7.

\subsection{Hệ thống AI hỗ trợ chấm điểm và cá nhân hóa luyện tập kỹ năng Nói IELTS}
\label{subsec:wispace_system_spec}

\subsubsection{Mô tả tổng quan về hệ thống}
Wispace Speaking hỗ trợ luyện tập theo Part và thi thử IELTS với kho đề Cambridge 1--20, Forecast hoặc đề tự thêm. Scoring Model do Wispace tự huấn luyện nhận audio, trả transcript, lỗi phát âm và band theo bốn tiêu chí Fluency \& Coherence, Lexical Resource, Grammatical Range \& Accuracy và Pronunciation. Kết quả được dùng để theo dõi tiến bộ và xây dựng lộ trình cá nhân hóa.

Thi thử mô phỏng buổi thi 11--14 phút, có giọng đọc câu hỏi và không cho nghe lại hoặc ghi lại. Tài khoản, gói Self-practice và thanh toán PayOS được dùng chung trong nền tảng Wispace Writing/Speaking. Dữ liệu dùng để cải thiện mô hình cần sự đồng ý của học viên và nhãn đánh giá của chuyên gia.

\subsubsection{Các nhóm người dùng và chức năng}
\begin{itemize}
\item \textbf{Khách:} xem kho đề, đăng ký, đăng nhập và khôi phục mật khẩu.
\item \textbf{Học viên:} luyện tập, thi thử, xem kết quả, quản lý hồ sơ và mua gói dịch vụ.
\item \textbf{Kiểm duyệt viên:} quản lý kho đề, duyệt đề cộng đồng và xử lý báo cáo nội dung.
\item \textbf{Quản trị viên:} quản lý người dùng, gói, hạn mức và thống kê vận hành.
\end{itemize}
Hệ thống tích hợp Scoring Model, TTS, Google OAuth, email, PayOS và Tracking-System. PostgreSQL và S3 lưu dữ liệu nghiệp vụ và audio.

\subsubsection{Yêu cầu chức năng}
Các yêu cầu chức năng của hệ thống được phân rã chi tiết theo từng nhóm người dùng thông qua các bảng dưới đây.

\begin{table}[H]
\centering
\caption{Yêu cầu chức năng của Khách (chưa đăng nhập)}
\label{wispace:tab:func_guest}
\small
\begin{tabular}{|p{4.2cm}|p{\dimexpr\linewidth-4.2cm-4\tabcolsep-3\arrayrulewidth\relax}|}
\hline
\textbf{Chức năng} & \textbf{Mô tả} \\
\hline
Đăng ký tài khoản & Khách nhập email, mật khẩu. Hệ thống gửi mã OTP 6 chữ số qua email để kích hoạt tài khoản, sau đó cấp gói dùng thử và chuyển người dùng sang bước onboarding. \\
\hline
Đăng nhập & Đăng nhập bằng email và mật khẩu hoặc bằng tài khoản Google. Sau khi xác thực thành công, hệ thống chuyển người dùng đến trang tương ứng với vai trò của họ. \\
\hline
Quên mật khẩu & Khách nhập email đã đăng ký, nhận mã OTP qua email, xác thực OTP rồi đặt mật khẩu mới. \\
\hline
Xem danh sách đề & Xem danh sách đề thi Cambridge và Forecast (chỉ hiển thị tiêu đề và chủ đề). Khi bấm luyện tập, hệ thống yêu cầu đăng nhập. \\
\hline
\end{tabular}
\end{table}

\begin{table}[H]
\centering
\caption{Yêu cầu chức năng của Học viên}
\label{wispace:tab:func_student}
\small
\begin{tabular}{|p{4.2cm}|p{\dimexpr\linewidth-4.2cm-4\tabcolsep-3\arrayrulewidth\relax}|}
\hline
\textbf{Chức năng} & \textbf{Mô tả} \\
\hline
Đăng xuất & Đăng xuất khỏi thiết bị hiện tại hoặc tất cả thiết bị. Nếu đang làm bài thi thử, hệ thống hỏi người dùng có muốn nộp phần đã làm hay không. Hệ thống vô hiệu hoá token và đưa người dùng về trang đăng nhập. \\
\hline
Quản lý hồ sơ cá nhân & Cập nhật tên, ảnh đại diện, band mục tiêu, ngày thi dự kiến và giọng giám khảo ưa thích. Kiểm tra micro, xem gói đang dùng và hạn mức còn lại, đổi mật khẩu, xoá file ghi âm hoặc xoá tài khoản. \\
\hline
Xem và tìm kiếm đề thi & Duyệt đề theo các nguồn: Cambridge 1--20, Forecast theo quý, Đề của tôi và Cộng đồng. Lọc theo Part, chủ đề và trạng thái đã làm. Từ trang chi tiết đề, người dùng chọn Luyện tập hoặc Thi thử. \\
\hline
Thêm đề thi & Tự tạo đề cho Part 1, Part 2 (cue card), Part 3 hoặc cả bộ đề. Hệ thống kiểm tra tính hợp lệ và phát hiện trùng lặp, lưu đề ở chế độ riêng tư; người dùng có thể đề xuất công khai để Moderator duyệt. \\
\hline
Luyện tập theo Part & Chọn câu hỏi, nghe giọng đọc (TTS) của giám khảo, ghi âm câu trả lời, nghe lại hoặc ghi âm lại trước khi nộp. Với Part 2 có 60 giây chuẩn bị kèm giấy nháp và thời gian nói tối đa 2 phút. Audio được gửi cho Scoring Model chấm và trả kết quả từng câu. \\
\hline
Xem kết quả và luyện lại & Xem band tổng và 4 tiêu chí (FC, LR, GRA, P), transcript có tô lỗi, chỉ số độ trôi chảy và câu trả lời mẫu. Bấm vào từ phát âm sai để nghe lại và so sánh với phát âm chuẩn. Có thể luyện lại và so sánh kết quả giữa lần trước và lần này. \\
\hline
Thi thử như thật & Chọn bộ đề đầy đủ và kiểm tra thiết bị. Giám khảo AI hỏi bằng giọng nói lần lượt Part 1, Part 2, Part 3. Không cho nghe lại, ghi âm lại hay tạm dừng. Khi kết thúc, hệ thống chấm bài ở chế độ nền và thông báo khi báo cáo sẵn sàng. \\
\hline
Xem lịch sử và tiến bộ & Xem danh sách các bài đã làm, lọc theo Part, nguồn đề và ngày làm. Theo dõi biểu đồ band theo thời gian và danh sách các lỗi lặp lại nhiều nhất. \\
\hline
Mua gói dịch vụ & Chọn gói Self-practice (1, 3, 6 hoặc 12 tháng), dùng chung cho Writing và Speaking. Hệ thống tạo liên kết thanh toán PayOS; sau khi PayOS xác nhận, quyền lợi trả phí được kích hoạt cho cả Writing và Speaking. \\
\hline
\end{tabular}
\end{table}

\begin{table}[H]
\centering
\caption{Yêu cầu chức năng của Moderator}
\label{wispace:tab:func_moderator}
\small
\begin{tabular}{|p{4.2cm}|p{\dimexpr\linewidth-4.2cm-4\tabcolsep-3\arrayrulewidth\relax}|}
\hline
\textbf{Chức năng} & \textbf{Mô tả} \\
\hline
Quản lý kho đề Cambridge/Forecast & Thêm, sửa, ẩn đề theo cấu trúc Quyển $\rightarrow$ Test $\rightarrow$ Part $\rightarrow$ Câu hỏi. Nhập hàng loạt bằng file CSV/JSON và nhận báo cáo các dòng lỗi. Tạo bộ Forecast theo quý và sinh audio cho câu hỏi. \\
\hline
Duyệt đề cộng đồng & Xem hàng chờ duyệt, sau đó duyệt, từ chối kèm lý do hoặc gộp với đề trùng. Hệ thống thông báo kết quả cho tác giả. \\
\hline
Xử lý báo cáo đề & Xem các báo cáo đề sai, trùng hoặc vi phạm. Moderator quyết định gỡ đề, sửa đề hoặc bỏ qua báo cáo. \\
\hline
\end{tabular}
\end{table}

\begin{table}[H]
\centering
\caption{Yêu cầu chức năng của Quản trị viên hệ thống}
\label{wispace:tab:func_admin}
\small
\begin{tabular}{|p{4.2cm}|p{\dimexpr\linewidth-4.2cm-4\tabcolsep-3\arrayrulewidth\relax}|}
\hline
\textbf{Chức năng} & \textbf{Mô tả} \\
\hline
Tìm kiếm người dùng & Tìm theo email hoặc tên; lọc theo vai trò, gói dịch vụ, trạng thái và ngày đăng ký. Xem trang chi tiết của từng người dùng. \\
\hline
Cập nhật trạng thái người dùng & Khoá hoặc mở khoá tài khoản, đổi vai trò, cộng thêm quota thủ công. Bắt buộc nhập lý do; thay đổi được áp dụng ngay, ghi nhật ký và gửi email thông báo cho người dùng. \\
\hline
Quản lý gói và quota & Cấu hình gói dùng chung, giá và hạn mức từng tính năng cho Writing và Speaking: số bài viết được chấm, số câu nói được chấm và số lượt thi thử mỗi ngày. \\
\hline
Xem thống kê & Xem doanh thu, số người dùng hoạt động, số bài Speaking đã chấm và chi phí AI/STT theo ngày. \\
\hline
\end{tabular}
\end{table}
\newpage

\subsubsection{Yêu cầu phi chức năng}
Các yêu cầu dưới đây xác định chất lượng phục vụ, khả năng bảo vệ dữ liệu và tính ổn định của Wispace Speaking. Các giá trị định lượng là mục tiêu kiểm thử khi hiện thực hệ thống.

\begin{itemize}
\item \textbf{Hiệu năng và khả năng chịu tải:} kết quả chấm một câu có audio không quá 60 giây được trả trong tối đa 15 giây; báo cáo thi thử khoảng 12 phút audio sẵn sàng trong tối đa 3 phút sau khi kết thúc. Câu hỏi tiếp theo được phát trong tối đa 2 giây và trang kho 2.000 đề tải trong tối đa 2 giây. Hệ thống hướng đến 1.000--10.000 người dùng đồng thời; tác vụ chấm nền không cản trở thao tác ghi âm hoặc chuyển câu.

\item \textbf{Chất lượng chấm điểm:} ít nhất 80\% bài trong bộ mẫu chuẩn có band lệch không quá 0.5 so với giám khảo; WER không quá 15\% đối với giọng người Việt. Transcript giữ nguyên lỗi của người nói để phục vụ đánh giá. Kết quả thể hiện bốn tiêu chí FC, LR, GRA, P, phiên bản mô hình và thông báo ``Điểm tham khảo do AI chấm''; bài không đủ dữ liệu phải được phân biệt với bài đã chấm đầy đủ.

\item \textbf{Thu âm và truyền dữ liệu:} hỗ trợ WebM/Opus và MP4/AAC trên Safari, chuẩn hóa về 16 kHz mono khi xử lý. Audio được tải theo đoạn 5--10 giây; gián đoạn mạng dưới 60 giây không làm mất phần ghi âm. Hệ thống cảnh báo mic nhỏ, rè, ồn hoặc im lặng quá 80\% đoạn ghi; các đoạn được xác nhận đầy đủ trước khi nộp bài.

\item \textbf{Bảo mật và phân quyền:} mật khẩu được băm, giao tiếp sử dụng HTTPS và các API kiểm tra vai trò, quyền sở hữu bài làm. Audio được mã hóa, lưu riêng tư và chỉ truy cập bằng liên kết ký có thời hạn tối đa 10 phút. Khóa đăng nhập 15 phút sau 5 lần sai; giới hạn 30 phiên ghi âm/giờ/người. Nội dung đề tự thêm được kiểm tra để loại bỏ mã độc và dữ liệu gây can thiệp vào chỉ dẫn AI.

\item \textbf{Quyền riêng tư và vòng đời dữ liệu:} học viên xác nhận đồng ý xử lý giọng nói khi sử dụng lần đầu. Audio được xóa sau 180 ngày nhưng giữ transcript và điểm theo chính sách; người dùng có thể yêu cầu xóa audio hoặc tài khoản. Sự đồng ý dùng dữ liệu để huấn luyện được quản lý riêng, không mặc định đi kèm việc sử dụng chức năng chấm điểm.

\item \textbf{Tính sẵn sàng và khôi phục:} mục tiêu hoạt động tối thiểu 99.5\%/tháng. Bài đã nộp vẫn được giữ khi dịch vụ AI gián đoạn; mỗi công việc có tối đa ba lần thực thi tự động, gồm lần đầu và hai lần thử lại. Thất bại cuối cùng được hoàn hạn mức theo nghiệp vụ, không hoàn hoặc trừ lặp khi gửi lại yêu cầu. Phiên thi có thể khôi phục từ trạng thái đã lưu; TTS có phương án dự phòng.

\item \textbf{Khả năng sử dụng và tương thích:} hỗ trợ Chrome, Edge, Safari, giao diện responsive tiếng Việt/Anh và phím Space để bắt đầu hoặc dừng ghi âm. Trạng thái ghi âm, tải lên và chấm bài được hiển thị rõ; lỗi cần có thông báo và hướng xử lý phù hợp. Phòng thi thử khuyến nghị máy tính và tai nghe, có bước kiểm tra thiết bị trước khi bắt đầu.

\item \textbf{Bảo trì và giám sát:} kết quả gắn với phiên bản mô hình để truy vết và so sánh chất lượng. Hệ thống ghi mã truy vết theo phiên, thời gian xử lý và chi phí AI/STT/TTS theo người dùng, bài làm; cảnh báo khi vượt ngưỡng ngày. Audio câu hỏi cố định được tạo một lần và tái sử dụng. Nhật ký quản trị ghi người thực hiện, thời điểm và lý do thay đổi.
\end{itemize}

\subsubsection{Yêu cầu dữ liệu}
\label{wispace:subsubsec:wispace_data_requirements}
Dữ liệu Wispace cần mô tả được toàn bộ quá trình từ lựa chọn đề, ghi âm, nộp bài đến chấm điểm và theo dõi tiến bộ. Các nhóm dữ liệu chính gồm:

\begin{itemize}
\item \textbf{Tài khoản, hồ sơ và phiên truy cập:} email, mật khẩu băm, liên kết Google, vai trò, trạng thái tài khoản, OTP và phiên xác thực. Hồ sơ học viên gồm tên, ảnh đại diện, band mục tiêu, ngày thi dự kiến và giọng đọc câu hỏi. Các thay đổi mật khẩu, đăng xuất và khóa tài khoản phải liên kết được với trạng thái phiên.

\item \textbf{Gói dịch vụ, hạn mức và thanh toán:} gói Self-practice 1/3/6/12 tháng dùng chung Writing/Speaking, thời hạn và quyền lợi từng tính năng. Lưu lượng sử dụng, khoản giữ, trừ và hoàn hạn mức phải gắn với bài làm hoặc phiên thi cụ thể. Giao dịch PayOS lưu mã đơn, số tiền và trạng thái xác nhận để tránh kích hoạt quyền lợi nhiều lần.

\item \textbf{Kho đề và kiểm duyệt nội dung:} nguồn Cambridge/Forecast/đề cá nhân, Test, Part, câu hỏi, chủ đề, cue card và audio câu hỏi. Dữ liệu đề có tác giả, phiên bản, phạm vi công khai và trạng thái kiểm duyệt; đề xuất cộng đồng kèm quyết định duyệt, lý do từ chối hoặc báo cáo vi phạm. Câu hỏi dùng trong phiên thi cần được bảo toàn khi kho đề thay đổi.

\item \textbf{Phiên luyện tập, thi thử và audio:} học viên, đề đã chọn, thời điểm bắt đầu/kết thúc, thời hạn từng phần, câu trả lời và trạng thái phiên. Metadata audio gồm định danh tệp, vị trí lưu trên S3, định dạng, thời lượng, kích thước, checksum và thứ tự các đoạn. Dữ liệu thể hiện được đoạn nào đã tải, đã xác minh hoặc còn thiếu để tiếp tục khi kết nối gián đoạn.

\item \textbf{Công việc chấm và kết quả:} mã bài nộp, trạng thái chấm, lịch sử thực thi, lỗi xử lý và phiên bản mô hình. Kết quả gồm transcript có mốc thời gian, band tổng, điểm FC/LR/GRA/P, nhận xét, lỗi phát âm/ngữ pháp/từ vựng và chỉ số trôi chảy. Bài được chấm lại phải giữ liên kết với lần trước, không ghi đè mất lịch sử.

\item \textbf{Tiến bộ và lộ trình cá nhân:} lịch sử bài luyện/thi thử, diễn biến điểm, lỗi lặp và mục tiêu cải thiện. Các thống kê cần truy về bài làm và khoảng thời gian tổng hợp. Liên kết chia sẻ chỉ cung cấp phạm vi điểm/nhận xét được cho phép, không mặc định công khai audio và có thể thu hồi.

\item \textbf{Quyền riêng tư và vận hành:} bản ghi đồng ý xử lý giọng nói, sự đồng ý huấn luyện, yêu cầu xóa và trạng thái thực hiện. Nhật ký quản trị lưu người thao tác, thời điểm, giá trị thay đổi và lý do; dữ liệu vận hành lưu độ trễ, lỗi và chi phí dịch vụ để đối soát.
\end{itemize}

Dữ liệu cần nhất quán giữa trạng thái bài làm, kết quả chấm và hạn mức; mỗi học viên chỉ truy cập dữ liệu thuộc quyền của mình. Metadata phải khớp với tệp thực tế trên S3. Audio tuân theo thời hạn lưu trữ; yêu cầu xóa tài khoản có thời gian chờ 30 ngày. Dự đoán AI và nhãn chuẩn của giám khảo được phân biệt để phục vụ đánh giá và cải thiện mô hình.

\subsection{Hệ thống cung cấp API và quản trị hệ thống - AIHub}
\label{subsec:aihub_system_spec}

\subsubsection{Mô tả tổng quan về hệ thống}
AIHub cung cấp và quản lý API mô hình cho các tổ chức B2B. Developer Portal hỗ trợ quản lý khóa, kiểm thử dịch vụ, theo dõi hạn mức và nhật ký. Gateway xác thực, kiểm soát quyền và điều phối yêu cầu đến dịch vụ AI.

Academic Lab là công cụ nội bộ để sinh, hiệu đính, chấm chất lượng và thẩm định mẫu học thuật. Bản gốc AI, bản hiệu đính và phản hồi được lưu làm dữ liệu phục vụ SFT/DPO. Wispace Speaking có đường tích hợp mô hình riêng, không phụ thuộc cổng B2B của AIHub.

\subsubsection{Các nhóm người dùng và chức năng}
\begin{itemize}
\item \textbf{Đối tác tích hợp:} quản lý khóa, kiểm thử API và tra cứu usage của tổ chức.
\item \textbf{Chuyên viên học thuật:} tra cứu dẫn chứng, sinh mẫu, hiệu đính và thẩm định nội dung.
\item \textbf{Kỹ sư AI:} quản lý Prompt, cấu hình mô hình, phân tích lỗi và xuất dữ liệu.
\item \textbf{Quản trị viên tổ chức:} quản lý thành viên, khóa và hạn mức trong tổ chức.
\item \textbf{Quản trị viên nền tảng:} quản lý danh mục dịch vụ, gói, giám sát và kiểm toán.
\end{itemize}
Các dịch vụ tích hợp gồm mô hình AI, nhà cung cấp LLM và API dẫn chứng Wispace được cấp quyền.

\subsubsection{Yêu cầu chức năng}
Các yêu cầu chức năng của hệ thống được phân rã chi tiết theo từng nhóm người dùng thông qua các bảng dưới đây.

\begin{table}[H]
\centering
\caption{Danh sách các chức năng của Lập trình viên / Đối tác tích hợp}
\label{aihub:tab:func_developer}
\small
\begin{tabular}{|p{4.2cm}|p{\dimexpr\linewidth-4.2cm-4\tabcolsep-3\arrayrulewidth\relax}|}
\hline
\textbf{Chức năng} & \textbf{Mô tả} \\ \hline
Đăng nhập \& Xác thực & Lập trình viên đăng nhập thông qua phiên Web/JWT hoặc sử dụng mã khóa \textit{x-client-id} và \textit{x-secret-key} khi gọi trực tiếp qua giao diện API. \\ \hline
Quản lý khóa API (Credentials) & Tạo mới, liệt kê danh sách, tái tạo (rotate) và vô hiệu hóa (revoke) cặp mã định danh \textit{ClientId} và bí mật \textit{SecretKey}. \\ \hline
Kiểm thử cURL Sandbox & Lựa chọn đặc tả dịch vụ trong danh mục Universal Specs, tùy biến tham số đầu vào và thực thi kiểm thử trực tiếp trên giao diện để kiểm tra tính hợp lệ của lược đồ JSON trả về. \\ \hline
Gọi API sinh bài tập vi mô & Tích hợp gọi endpoint \texttt{/create-micro-exercise} hoặc \texttt{/ai/v1/create-micro-exercise} với các tham số mã lỗi (28 mã lỗi IELTS), dạng bài và nhận cấu trúc bài tập hoàn chỉnh. \\ \hline
Tra cứu hạn mức \& Nhật ký sử dụng & Theo dõi lượng định ngạch tiêu thụ theo tháng (Monthly Quota), giới hạn tần suất yêu cầu trên phút (RPM) và lịch sử gọi API chi tiết (thời gian, mã trạng thái, số token tiêu thụ). \\ \hline
\end{tabular}
\end{table}

\begin{table}[H]
\centering
\caption{Danh sách các chức năng của Chuyên viên học thuật (Academic Lab)}
\label{aihub:tab:func_academic}
\small
\begin{tabular}{|p{4.2cm}|p{\dimexpr\linewidth-4.2cm-4\tabcolsep-3\arrayrulewidth\relax}|}
\hline
\textbf{Chức năng} & \textbf{Mô tả} \\ \hline
Truy vấn ngân hàng lỗi thực tế & Tra cứu ngữ cảnh lỗi viết thực tế của học viên từ cơ sở dữ liệu hoặc hệ thống WISPACE theo Task (1 hoặc 2), tiêu chí đánh giá (TR, CC, LR, GRA) và chủ đề bài học. \\ \hline
Khởi tạo mẫu bài tập học thuật & Gửi yêu cầu sinh mẫu bài tập mới dựa trên bộ tham số sư phạm cụ thể; hệ thống ghi nhận bản kết quả AI thô (\textit{rawAiOutput}) vào cơ sở dữ liệu. \\ \hline
Kiểm định \& Chỉnh sửa bài tập & Xem chi tiết bài tập được tạo, kiểm tra tính sư phạm, trực tiếp chỉnh sửa câu hỏi, đáp án, giải thích và tiêu chí chấm điểm để tạo phiên bản hiệu đính (\textit{finalOutput}). \\ \hline
Đánh giá thang điểm chất lượng & Chấm điểm chất lượng bài tập trên thang điểm từ 1 đến 10 theo các tiêu chuẩn về độ chính xác ngữ pháp, tính tự nhiên và độ bám sát chuẩn IELTS. \\ \hline
Chấp thuận mẫu (Accept) & Xác nhận lưu trữ đồng thời cặp dữ liệu \textit{rawAiOutput} và \textit{finalOutput} cùng điểm số vào cơ sở dữ liệu ở trạng thái ACCEPTED; hệ thống tự động gợi ý bước xoay vòng chủ đề và mã lỗi tiếp theo. \\ \hline
Từ chối mẫu (Reject) & Ghi nhận mẫu ở trạng thái REJECTED kèm lý do từ chối cụ thể (\textit{rejectReason}), chuyển dữ liệu sang phân hệ xử lý mẫu âm cho mô hình học máy. \\ \hline
Quản lý kho dữ liệu mẫu & Xem danh sách kiểm kê toàn bộ các mẫu bài tập đã sinh, hỗ trợ lọc theo trạng thái (DRAFT, ACCEPTED, REJECTED), mã lỗi, người thẩm định và đánh dấu mục yêu thích. \\ \hline
\end{tabular}
\end{table}

\begin{table}[H]
\centering
\caption{Danh sách các chức năng của Kỹ sư AI}
\label{aihub:tab:func_ai_engineer}
\small
\begin{tabular}{|p{4.2cm}|p{\dimexpr\linewidth-4.2cm-4\tabcolsep-3\arrayrulewidth\relax}|}
\hline
\textbf{Chức năng} & \textbf{Mô tả} \\ \hline
Cấu hình mẫu Prompt & Quản lý nội dung hướng dẫn hệ thống, cấu trúc ràng buộc định dạng JSON cho từng loại mã lỗi và từng dạng câu hỏi (trắc nghiệm, điền từ, viết ngắn). \\ \hline
Cấu hình chuỗi mô hình dự phòng & Thiết lập thứ tự ưu tiên các nhà cung cấp mô hình (WISPACE Proxy $\rightarrow$ Google Gemini $\rightarrow$ OpenAI GPT-4o $\rightarrow$ Deterministic Engine). \\ \hline
Trích xuất tập dữ liệu tinh chỉnh & Xuất các cặp dữ liệu đã qua thẩm định học thuật (\textit{rawAiOutput}, \textit{finalOutput}, điểm chất lượng và lý do từ chối) thành định dạng phục vụ huấn luyện mô hình SFT/DPO. \\ \hline
Theo dõi phân tích lỗi mô hình & Phân tích tần suất các mã lỗi bị từ chối, nhận diện hiện tượng suy thoái chất lượng sinh nội dung hoặc sai lệch lược đồ cấu trúc đầu ra. \\ \hline
\end{tabular}
\end{table}

\begin{table}[H]
\centering
\caption{Danh sách các chức năng của Quản trị viên (Tổ chức \& Nền tảng)}
\label{aihub:tab:func_admin}
\small
\begin{tabular}{|p{4.2cm}|p{\dimexpr\linewidth-4.2cm-4\tabcolsep-3\arrayrulewidth\relax}|}
\hline
\textbf{Chức năng} & \textbf{Mô tả} \\ \hline
Quản lý thành viên tổ chức & Tiếp nhận, cấp quyền, điều chỉnh vai trò thành viên (OWNER, ADMIN, MEMBER) trong phạm vi tổ chức thuê bao. \\ \hline
Quản trị gói cước \& Hạn mức & Thiết lập danh mục các gói thuê bao (Free Sandbox, Pro, Enterprise), quy định hạn ngạch yêu cầu hàng tháng và trần tần suất truy cập RPM tương ứng. \\ \hline
Giám sát trạng thái toàn hệ thống & Theo dõi endpoint thăm dò sức khỏe \texttt{/health}, trạng thái sẵn sàng của các nhà cung cấp LLM, độ trễ mạng trung bình và phân bố mã lỗi HTTP. \\ \hline
Quản lý danh mục Universal Specs & Cập nhật, kích hoạt hoặc điều chỉnh lược đồ JSON Schema của các dịch vụ AI liên miền trong cơ sở dữ liệu đặc tả hệ thống. \\ \hline
Kiểm toán hệ thống (Audit Logs) & Tra cứu tập trung nhật ký truy vết mọi giao dịch phát sinh qua Gateway nhằm phục vụ mục đích an ninh và tính cước doanh nghiệp. \\ \hline
\end{tabular}
\end{table}
\newpage

\subsubsection{Yêu cầu phi chức năng}
Các yêu cầu tập trung vào chất lượng cổng B2B, an toàn dữ liệu tổ chức và khả năng kiểm soát đầu ra AI. Các chỉ tiêu thời gian của tác vụ sinh nội dung không được áp dụng mặc định cho tác vụ chấm audio.

\begin{itemize}
\item \textbf{Hiệu năng và khả năng mở rộng:} chi phí xử lý tại Gateway không quá 50 ms. Sinh bài tập có dẫn chứng hoàn tất dưới 2.5 giây với mô hình chính, dưới 1.5 giây với động cơ xác định; Sandbox mô phỏng phản hồi trong 300--800 ms. Cổng B2B hướng đến tối thiểu 1.000 yêu cầu đồng thời; tác vụ dài có trạng thái để đối tác theo dõi mà không giữ kết nối đến khi hoàn tất.

\item \textbf{Xác thực và bảo vệ khóa:} bí mật API được sinh ngẫu nhiên, chỉ hiển thị một lần và lưu bằng Bcrypt; khóa có tiền tố phân biệt môi trường. Phiên Portal và credentials ứng dụng được quản lý riêng. Khóa đã thu hồi không được tiếp tục sử dụng; dữ liệu bí mật không xuất hiện trong log, phản hồi công khai hoặc bộ dữ liệu huấn luyện.

\item \textbf{Phân quyền và cô lập tổ chức:} RBAC giới hạn thao tác theo vai trò và tổ chức. Đối tác chỉ xem khóa, yêu cầu, hạn mức và usage của tổ chức mình; Academic Lab được cấp quyền nội bộ riêng. Hệ thống kiểm soát giới hạn tốc độ và hạn mức để một tổ chức không chiếm dụng toàn bộ năng lực dịch vụ.

\item \textbf{Tính toàn vẹn giao dịch:} yêu cầu gửi lại phải được nhận diện để tránh tạo công việc hoặc tính hạn mức trùng. Mức sử dụng cần liên kết với kết quả thực thi; khoản giữ hạn mức phải được chốt hoặc giải phóng theo trạng thái cuối. Chi phí suy luận phát sinh được phân biệt với số lượng hạn mức tính cho khách hàng để phục vụ đối soát.

\item \textbf{Chất lượng và tính nhất quán đầu ra:} dữ liệu vào/ra được kiểm tra theo schema của dịch vụ; request giới hạn 2 MB. Nội dung học viên được tách khỏi chỉ dẫn Prompt. Mô hình dự phòng phải tuân theo cùng hợp đồng, phản hồi ghi được mô hình thực tế sử dụng; khi không có kết quả hợp lệ, hệ thống trả lỗi có kiểm soát thay vì báo thành công với nội dung thiếu.

\item \textbf{Tính sẵn sàng và chịu lỗi:} mục tiêu hoạt động tối thiểu 99.9\%/tháng. Các dịch vụ sinh nội dung có chuỗi dự phòng và thời hạn xử lý theo cấu hình. Sự cố một nhà cung cấp không làm mất yêu cầu đã ghi nhận; công cụ giám sát phân biệt trạng thái bình thường, suy giảm và không sẵn sàng. Dịch vụ chấm Speaking chỉ sử dụng mô hình chấm phù hợp, không thay bằng LLM sinh nội dung.

\item \textbf{Bảo trì và truy vết:} nghiệp vụ được tách khỏi adapter nhà cung cấp; phiên bản Prompt, mô hình và đặc tả được quản lý để tái hiện kết quả. Nhật ký ghi tổ chức, dịch vụ, mã yêu cầu, thời gian, trạng thái và mức tiêu thụ. API hỗ trợ mở rộng theo chiều ngang, cấu hình thay đổi được kiểm tra và ghi nhận người thực hiện.
\end{itemize}

\subsubsection{Yêu cầu dữ liệu}
\label{aihub:subsubsec:aihub_data_requirements}
AIHub cần liên kết được danh tính tổ chức, dịch vụ khai thác, kết quả thực thi và mức sử dụng, đồng thời quản lý riêng dữ liệu Academic Lab. Các nhóm dữ liệu gồm:

\begin{itemize}
\item \textbf{Tài khoản, tổ chức và thành viên:} thông tin người dùng Portal, tổ chức, lời mời, vai trò OWNER/ADMIN/MEMBER, quyền nội bộ và trạng thái thành viên. Dữ liệu cho phép xác định ai được quản lý khóa, xem usage hoặc thực hiện thao tác học thuật; thay đổi quyền và chuyển quyền sở hữu được ghi nhận.

\item \textbf{Credentials và quyền khai thác dịch vụ:} clientId, mã băm bí mật, môi trường, tổ chức sở hữu, phạm vi quyền, thời điểm tạo và trạng thái thu hồi. Lịch sử tạo/xoay vòng/thu hồi phục vụ kiểm toán; không lưu lại bí mật gốc để cung cấp cho người dùng sau này.

\item \textbf{Danh mục API, gói và hạn mức:} mã dịch vụ, mô tả, phiên bản, schema yêu cầu/phản hồi, hình thức xử lý và đơn vị tính mức sử dụng. Gói xác định dịch vụ được phép dùng, hạn mức tháng, RPM và cấu hình riêng của tổ chức. Model route lưu mô hình, phiên bản và thứ tự dự phòng được phép cho từng dịch vụ.

\item \textbf{Yêu cầu API, usage và kiểm toán:} mã yêu cầu, tổ chức, credentials, dịch vụ, khóa chống trùng, thời điểm, trạng thái và kết quả hoặc lỗi. Usage lưu latency, mô hình thực tế, token hoặc thời lượng audio theo loại dịch vụ, cùng thông tin giữ/chốt/hoàn hạn mức. Nhật ký quản trị gắn người thực hiện và nội dung thay đổi để truy vết, giải quyết sai lệch.

\item \textbf{Ngữ cảnh và cấu hình học thuật:} đề Writing Task 1/2, tiêu chí TR/CC/LR/GRA, bộ 28 mã lỗi, chủ đề, trích dẫn và nhận xét chuyên môn. Prompt có phiên bản, biến đầu vào, schema và ví dụ mẫu; dẫn chứng cần xác định được nguồn và phạm vi sử dụng.

\item \textbf{Mẫu và lịch sử thẩm định:} giữ đồng thời rawAiOutput, finalOutput, phiên bản Prompt/mô hình và thông tin thẩm định. Điểm chất lượng là số nguyên 1--10; mẫu có trạng thái DRAFT, ACCEPTED hoặc REJECTED. Quyết định từ chối kèm lý do; lưu nháp và quyết định duyệt được phân biệt, có phiên bản để nhận diện xung đột chỉnh sửa.

\item \textbf{Bộ dữ liệu xuất:} định dạng SFT/DPO, bộ lọc, danh sách mẫu, phiên bản, người xuất, thời điểm, số mẫu và tệp kết quả. SFT liên kết Prompt với bản hiệu đính đạt tiêu chí; DPO giữ cặp đầu ra được chọn và bị từ chối cho cùng Prompt. Tệp xuất loại bỏ thông tin định danh cá nhân, có checksum và được lưu trong vùng nội bộ có kiểm soát truy cập.
\end{itemize}

Dữ liệu B2B được phân tách theo tổ chức, còn dữ liệu học thuật theo quyền nội bộ. Bản gốc AI không bị thay bằng bản hiệu đính; dự đoán của mô hình được phân biệt với nội dung đã thẩm định. Mỗi giao dịch, quyết định duyệt và lần xuất dữ liệu cần truy vết được nguồn gốc, trong khi các bí mật xác thực và dữ liệu cá nhân không thuộc nội dung xuất huấn luyện.

\subsection*{Kết chương}
Chương này đã xác định phạm vi nghiệp vụ, các tác nhân và các nhóm yêu cầu của AIHub và Wispace Speaking. Các nội dung này làm cơ sở cho đặc tả use-case và phân tích luồng hoạt động ở chương tiếp theo.
